\chapter{08 -- Pileup: the other 68 collisions of the bunch crossing, at generator level}
\label{ch:08_PU}

\section{What this chapter is about}

In 2025, the year of the release the book runs, a bunch crossing at CMS held on average some 60 inelastic
proton--proton collisions (64 in the 2025 profile of the release; the book takes 60). One of them is the \bbbar{} event the book follows; the others are \emph{pileup}, ordinary inelastic collisions
that happen to sit in the same 25\,ns window and that the detector cannot tell apart from the signal. Before
anything is simulated in the detector these collisions have to be generated, and this chapter does that
alone: it runs Pythia8 with the minimum-bias fragment of the release, \code{MinBias\_14TeV\_pythia8\_TuneCP5\_cfi},
at the $13.6\TeV$ of the signal, 250 times, and asks what an inelastic collision is (non-diffractive or
diffractive, how many parton pairs interact, how many hadrons come out, where the vertex is in space and
time). The 250 collisions are the generator-level \emph{pileup library} of the book; one of them is drawn
with the five generator pages of chapters 01--06.

The chapter then puts the pileup where it belongs, on top of the signal, still without a detector: the
stable particles of the 68 library collisions that will share the bunch crossing of our event in
chapters~\ref{ch:10_SIM} and \ref{ch:11_DIGI} are added to the 634 stable particles of chapter~\ref{ch:07_GenJet},
and the anti-$k_{\mathrm{T}}$ clustering is run again on the sum. The $b$ jet goes from 117.8 to
$150.9\GeV$ and 266 jets appear where there were 27: this is what pileup does to a jet before any detector
effect, and the number every pileup correction of the reconstruction will later try to undo.

\begin{plainwords}[Minimum bias]
The name is historical: a ``minimum-bias trigger'' was the loosest trigger an experiment could run, firing
on essentially any inelastic collision. A minimum-bias \emph{sample} is thus a sample of ordinary inelastic
collisions, with nothing selected and nothing enriched: the inelastic cross section itself, sampled in its
natural proportions, with no \ptHat{} cut and no filter. That is exactly what the other collisions of a bunch
crossing are, so a minimum-bias generator run is the raw material of pileup. In Pythia8 it is one switch,
\code{SoftQCD:inelastic = on}.
\end{plainwords}

\section{From the previous chapter}

Chapter~\ref{ch:07_GenJet} left the signal as a GEN file with its jets
(\path{test/07_GenJet/out/step2_GENJET_07_GenJet.root}): 1846 generator particles, 634 of them stable,
and 27 anti-$k_{\mathrm{T}}$ $R=0.4$ jets above $3\GeV$, the $b$ jet at $\pt=117.8\GeV$ ($\eta=-1.59$) and
the $\bar b$ jet at $115.9\GeV$ ($\eta=-3.75$). Nothing of it is rerun here (sample 8a); its jets alone are
the pages of chapter~\ref{ch:07_GenJet}, and \code{genJetAnalyzer\_08\_PU.pdf} holds the overlay.

The new ingredient is the library (sample 8b): 250 minimum-bias collisions generated by the GEN step of the
signal (vertex smearing from the database, energy) but with its own random seeds, so that its vertices are not
those of the signal. It is the first job of the book with the 2025 settings of the release, global tag
\code{auto:phase1\_2025\_realistic} and era \code{Run3\_2025}, as all the CMS chapters 08--13,
where the generator chapters 01--07 keep those of 2024; at generator level the two are the
same (section~\ref{sec:08_numbers}). The overlay (sample 8a+8b) takes library events 72 to 139 (\code{SKIP = 71}, \code{NPU = 68} in
\path{test/common/jitjet_common.sh}). Why these: the \code{MixingModule} of chapter~\ref{ch:10_SIM} (seed
12345, a Poisson draw with mean 60 per bunch crossing) reads the library in sequence, bunch crossing by bunch
crossing, early crossings first; for this signal event it draws 71 collisions for bunch crossing $-1$
(library events 1--71), 68 for bunch crossing 0 (events 72--139) and 54 for bunch crossing $+1$ (events
140--193). Library events 72--139 are therefore the in-time collisions of \emph{this} signal event, the same
ones, collision by collision, that the detector simulation will see in the same 25\,ns as the signal;
chapter~\ref{ch:10_SIM} checks its draw against them and chapter~\ref{ch:11_DIGI} replays it with the
playback of the \code{MixingModule}. Only in-time
collisions matter at generator level: a list of particles has no time axis, so the early and late crossings
can only act through the time response of the detector.

\section{The numbers of the event}
\label{sec:08_numbers}

\paragraph{The inelastic cross section.} Pythia8 prints, at the initialisation of the job, what it is
about to generate (table~\ref{tab:xsec}): ``We collide p+ with p+ at a CM energy of 1.360e+04 GeV'' and the
five subprocesses of \code{SoftQCD:inelastic} with their cross sections.\footnote{All numbers of this
section are from \code{test/08\_PU/out/step1\_GEN\_MinBias\_08\_PU.log} (the Pythia8 banner and statistics),
\code{gencheck\_08\_PU.log} and \code{genstats\_08\_PU.txt} (the catalogue of \code{common/minBiasGenStats.py}).}
The total and elastic cross sections are not computed: the tune CP5 fixes them
(\code{SigmaTotal:mode = 0}, \code{sigmaTot = 100.309}, \code{sigmaEl = 21.89} mb, in
\code{Configuration/Generator/python/MCTunes2017/PythiaCP5Settings\_cfi.py}), so that
$\sigma_{\mathrm{inel}} = 100.31 - 21.89 = 78.42\mb$ whatever the energy; the diffractive pieces come from
Pythia's own parametrisation and the non-diffractive cross section is the remainder. Pythia chooses the
process of every event in proportion to these numbers: 183 ND, 45 SD, 19 DD and 3 CD in the 250 events,
against 176.7, 41.0, 28.2 and 4.1 expected (the DD deficit is 1.7 standard deviations; the rest is within
one). The \code{GenXsecAnalyzer} at the end of the log reports the same $7.842\times10^{10}$\,pb with a filter
efficiency of $250/250$; the \code{crossSection = 71.39e+09} written in the fragment is old metadata that
nothing uses.

\begin{table}[htbp]
\centering\small\setlength{\tabcolsep}{4pt}
\caption{The inelastic collision in Pythia8 at $13.6\TeV$ with the tune CP5: the cross sections of the
``PYTHIA Process Initialization'' table of \code{step1\_GEN\_MinBias\_08\_PU.log}, and how many of the 250
collisions came out of each process (``Event and Cross Section Statistics'').}
\label{tab:xsec}
\begin{tabular}{llrrrr}
\toprule
process & code & $\sigma$ [mb] & share & expected of 250 & generated \\
\midrule
non-diffractive (ND), $pp\to X$        & 101       & 55.42 & 70.7\,\% & 176.7 & 183 \\
single diffractive (SD), $pp\to pX, Xp$ & 103 + 104 & $2\times6.43$ & 16.4\,\% & 41.0 & 45 (27 + 18) \\
double diffractive (DD), $pp\to XX$    & 105       & 8.85  & 11.3\,\% & 28.2 & 19 \\
central diffractive (CD), $pp\to pXp$  & 106       & 1.29  & 1.6\,\%  & 4.1 & 3 \\
\midrule
inelastic (\code{SoftQCD:inelastic})   & sum       & 78.42 & 100\,\%  & 250 & 250 \\
elastic, $pp\to pp$ (not generated)    & 102       & 21.89 & -- & -- & -- \\
total (fixed by CP5)                   &           & 100.31 & & & \\
\bottomrule
\end{tabular}
\end{table}

\begin{plainwords}[Four ways for two protons to collide inelastically]
Picture the two protons as bags of quarks and gluons. In a \emph{non-diffractive} collision the bags tear
open and their contents interact directly; the whole rapidity range fills with particles. In a
\emph{diffractive} collision the protons exchange something colourless (a ``Pomeron''), so that one of them
(SD) or both (DD) keep their identity or dissolve into a low-mass bunch of hadrons flying on along the beam,
and a gap with no particles separates the two sides. In a \emph{central-diffractive} collision both protons
survive and only a small system is created in the middle. The LHC does not pick: it makes all four in the
proportions of table~\ref{tab:xsec}, and so does the generator.
\end{plainwords}

\paragraph{Multiparton interactions set the size of a collision.} A non-diffractive collision has no hard
process to speak of; its multiplicity comes from the multiparton interactions of chapter~\ref{ch:04_MPI}.
The log shows the regularisation scale of the MPI cross section, $\pt^0 = 1.41\times(13.6/7)^{0.03344} =
1.44\GeV$ (\code{pT0Ref} and \code{ecmPow} of CP5), and ``pT0 = 1.44 gives sigmaInteraction = 492.18 mb'':
the integrated parton--parton cross section above the cut-off is nine times the non-diffractive cross
section of $55.42\mb$, so on average nine parton pairs interact per ND collision. The catalogue confirms it:
counting the hardest scattering plus one interaction per pair of status-31 partons, the 183 ND collisions
have between 1 and 39 interactions, mean 9.0, and the number of stable particles follows them almost
linearly, from 85 for one interaction (44 collisions) to 723 for 39 (the largest collision of the file,
event 247). The diffractive collisions have on average 1.2 to 2.0 interactions and are nearly empty for the central
detector: 70 stable particles and 7 central ones for SD, 57 and 2 for DD, 35 and 2 for CD, against 241 and
27 for ND (table~\ref{tab:cat}).

\begin{table}[htbp]
\centering\small
\caption{The 250 minimum-bias collisions at generator level (\code{genstats\_08\_PU.txt}): per collision,
the stable particles (HepMC status 1), the charged ones, the ``central'' ones ($\abseta<2.5$,
$\pt>0.5\GeV$), those beyond $\abseta=5.5$ where no CMS subdetector is, the hardest particle, the scalar
sum of \pt{}, the anti-$k_{\mathrm{T}}$ $R=0.4$ jets above $10\GeV$, and the smeared vertex. By process:
183 ND / 45 SD / 19 DD / 3 CD collisions with on average 241 / 70 / 57 / 35 stable and 27 / 7 / 2 / 2
central particles.}
\label{tab:cat}
\begin{tabular}{lrrrr}
\toprule
per collision & mean & median & min & max \\
\midrule
parton--parton interactions (1 + MPI) & 7.0 & 3 & 1 & 39 \\
stable particles & 193.8 & 132 & 3 & 723 \\
charged stable particles & 89.0 & 59 & 2 & 320 \\
central ($\abseta<2.5$, $\pt>0.5\GeV$) & 21.2 & 9.5 & 0 & 126 \\
beyond $\abseta=5.5$ & 38.9 & 36.5 & 3 & 93 \\
\pt{} of the hardest particle [GeV] & 2.20 & 1.81 & 0.23 & 8.08 \\
scalar sum of \pt{} [GeV] & 79.2 & 44.7 & 0.5 & 376.1 \\
jets above $10\GeV$ & 0.15 & 0 & 0 & 3 \\
vertex $z$ [cm] (r.m.s.\ 3.56) & 0.08 & $-0.05$ & $-9.5$ & 10.1 \\
vertex $t$ [ns] (r.m.s.\ 0.122) & $-0.007$ & $-0.015$ & $-0.39$ & 0.37 \\
\bottomrule
\end{tabular}
\end{table}

\paragraph{What a minimum-bias collision is made of.} Of the 48\,437 stable particles of the 250
collisions, 45.9\,\% are photons (almost all from $\pi^0\to\gamma\gamma$, decayed by Pythia because
$c\tau(\pi^0)$ is far below the 10\,mm of \code{ParticleDecays:tau0Max = 10}), 37.1\,\% charged pions,
4.8\,\% charged kaons, 2.7\,\% protons, 2.5\,\% neutrons, 2.4\,\% $K^0_L$, 2.2\,\% $K^0_S$ and 0.9\,\%
$\Lambda$ (all three left to Geant4). They are soft and forward: 36\,000 of them have $\pt<0.5\GeV$ and 17 have
$\pt>5\GeV$; the density per unit of $\eta$ is flat out to $\abseta=4$ and 20\,\% of the particles are
beyond $\abseta=5.5$. The signal event has two partons at $\pt=113\GeV$; a minimum-bias collision has
nothing above a few GeV and 0.15 jets above $10\GeV$.

\paragraph{The one collision drawn.} Figure~\ref{fig:mb} is library event 72, the first of the 68 in-time
collisions, at the end of the five generator pages of \code{genPartAnalyzer\_cfg.py}: 2 particles on the
hard-process page (a $2\to2$ scattering below $5\GeV$, not the $113\GeV$ of the signal), 10 after its
showers, 26 with the MPI, their showers and the beam remnants, 102 after hadronisation, 192 stable after the decays
(\code{genPartAnalyzer\_MinBias\_08\_PU.log}); a smaller than average ND collision, with 3 interactions. The
68 in-time collisions together (events 72--139 of the catalogue) are 49 ND, 15 SD and 4 DD with 467
interactions and 12\,993 stable particles, from 3 (event 132, SD) to 547 (event 94, ND), median 150.5: a sum
of very unequal terms.

\begin{figure}[htbp]
\centering
\includegraphics[width=0.8\textwidth,page=5]{08_PU/out/genPartAnalyzer_MinBias_08_PU.pdf}
\caption{One pileup collision at the end of the generator chain: library event 72, the first in-time
collision of the mixing, drawn like the display of chapter 06 (page 5 of \code{genPartAnalyzer\_MinBias\_08\_PU.pdf}).
A non-diffractive collision with 3 parton--parton interactions, 192 stable particles (84 charged, 13
central, 31 beyond $\abseta=5.5$, the 18 beyond $\abseta=6$ drawn on the border), the hardest at $2.2\GeV$, $\sum\pt = 56.6\GeV$,
vertex at $z=-0.48\cm$, $t=-0.092\ns$. No $b$, no $\bar b$, no hard process.}
\label{fig:mb}
\end{figure}

\paragraph{The same GEN step as the signal, with its own seeds.} The library is one \code{cmsDriver.py}
command (companion, section 1) with three lines appended: \code{process.generator.comEnergy = 13600} (the
fragment says $14\TeV$; the release has no $13.6\TeV$ CP5 minimum-bias fragment), \code{generator.initialSeed
= 20241} and \code{VtxSmeared.initialSeed = 20242}. The seeds matter: \code{cmsDriver.py} writes the same
default seeds into every configuration, so a library generated with the defaults would have, event by
event, the vertex of signal event $k$ in library event $k$, and the first pileup collision would sit exactly
on the signal vertex. The vertex smearing is the one of the signal, \code{VtxSmearedRealistic\_cfi}
(\code{BetafuncEvtVtxGenerator} reading the \code{SimBeamSpotObjectsRcd} payload of the global tag): a Gaussian in
$z$ and, with the same width in units of $c$, in time; the 250 vertices have r.m.s.\ $3.56\cm$ in $z$ and
$122\ps$ in $t$. The payload (\path{SimBeamSpot_Realistic25ns13p6TeVEarly2024CollisionVtxSmearingParameters_v1_mc},
centre $(0.028, -0.017, 0.42)\cm$, $\sigma_z = 3.6\cm$) was made for early 2024; the 2025 global tag keeps it. The time is part of the generated event: every later step starts the particles at it, and the mixing
only adds the bunch-crossing offset.

\paragraph{2024 or 2025: the same library.} The two global tags differ in 35 of about 620 records (companion,
section 4.9: calorimeter and pixel calibrations, alignment, offline beam spot, Level-1 menu), none read by a GEN
job, and \code{Run3\_2025} adds to \code{Run3\_2024} only modifiers of the reconstruction and the Level-1
trigger: rerun with the 2025 settings, the library and the overlay are the 2024 ones collision by collision.

\paragraph{The overlay.} No new code touches the particles (\path{python/genJetPileup_cff.py}): the
\code{MixingModule} of the DIGI step, with an empty \code{digitizers} set and \code{mixHepMC} as its only
mixed object, reads the \code{generatorSmeared} HepMC record of the signal and of 68 library events
(\code{sequential = True}, \code{skipEvents = 71}, \code{type = fixed}) into a
\code{CrossingFrame<HepMCProduct>} for the single bunch crossing \code{minBunch = maxBunch = 0}. The
\code{GenParticleProducer} of every GEN job, with \code{useCrossingFrame = True}, turns every event of the
frame into one collection, \code{genParticlesPU}, the signal first with \code{collisionId} 0 and the pileup
collisions as 1 to 68; the selector and the \code{ak4GenJets} producer of chapter~\ref{ch:07_GenJet} then
run on it as \code{ak4GenJetsPU}. The page is the same mixing as chapter~\ref{ch:11_DIGI}, stopped before
the detector.

\section{Ledger and display}

Table~\ref{tab:ledger} is the ledger of the overlay: the final state of the signal, of the 68 pileup
collisions and of the sum, by kind and by region, and the jets before and after.\footnote{From
\code{test/08\_PU/out/ledger\_08\_PU.txt} (\code{common/genLedger.py --pu}) and \code{genJetPileup\_08\_PU.log};
the signal column is the ledger of chapter~\ref{ch:07_GenJet}. The jet counts above 20 / 30 / $50\GeV$, the
dijet masses and the signal / pileup split of the constituents are computed from the jet lists of the two
logs and from \code{step2\_GENJETPU\_08\_PU.root}.} Pileup adds 12\,993 stable particles to the 634 of the
signal, $5315\GeV$ of scalar \pt{} to its 529, and $925\TeV$ of energy (68 collisions of $13.6\TeV$ each,
91\,\% of it beyond $\abseta=5.5$ in the beam remnants). In the tracker region it is $2179\GeV$ of \pt{}
spread over 4864 particles of $0.45\GeV$ each on average: a pedestal of $2179 / (5\times2\pi) = 69\GeV$ per
unit of $\eta$--$\phi$ area, flat within what one event can tell. A cone of $R=0.4$ covers $\pi R^2 = 0.50$
units and so collects about $35\GeV$ wherever it is put; the $b$ jet collected 33.9 and the $\bar b$ jet
$36.7\GeV$ (90 and 70 pileup constituents from 41 and 32 different collisions, on top of their 24 and 18
signal particles, which are the same as in chapter~\ref{ch:07_GenJet}).

\begin{table}[htbp]
\centering\footnotesize\setlength{\tabcolsep}{3.5pt}
\caption{The ledger of the overlay: the final state (status 1) of the signal, of the 68 in-time pileup
collisions and of the sum, by kind and by region, and the anti-$k_{\mathrm{T}}$ $R=0.4$ jets before and after
the overlay.}
\label{tab:ledger}
\begin{tabular}{lrrrrrrrrr}
\toprule
 & \multicolumn{3}{c}{signal} & \multicolumn{3}{c}{pileup (68 collisions)} & \multicolumn{3}{c}{signal + pileup} \\
\cmidrule(lr){2-4}\cmidrule(lr){5-7}\cmidrule(lr){8-10}
final state & $n$ & $E$ [GeV] & \pt{} [GeV] & $n$ & $E$ [GeV] & \pt{} [GeV] & $n$ & $E$ [GeV] & \pt{} [GeV] \\
\midrule
charged hadrons & 303 & 7168 & 259.1 & 5911 & 572\,735 & 3087.8 & 6214 & 579\,903 & 3347.0 \\
neutral hadrons & 56 & 3450 & 84.7 & 1088 & 179\,328 & 749.7 & 1144 & 182\,778 & 834.4 \\
photons & 273 & 2876 & 142.7 & 5900 & 171\,676 & 1458.8 & 6173 & 174\,552 & 1601.4 \\
electrons, muons & 1 & 11 & 5.0 & 87 & 1004 & 12.6 & 88 & 1015 & 17.6 \\
neutrinos & 1 & 96 & 37.5 & 7 & 58 & 5.6 & 8 & 153 & 43.0 \\
\midrule
$\abseta<2.5$ & 316 & 734 & 278.7 & 4864 & 5519 & 2179.4 & 5180 & 6253 & 2458.1 \\
$2.5<\abseta<3.0$ & 55 & 193 & 24.6 & 1037 & 3597 & 449.2 & 1092 & 3790 & 473.9 \\
$3.0<\abseta<5.2$ & 189 & 5020 & 197.4 & 3996 & 58\,413 & 1674.3 & 4185 & 63\,433 & 1871.7 \\
$5.2<\abseta<5.5$ & 16 & 836 & 7.8 & 436 & 17\,054 & 160.3 & 452 & 17\,890 & 168.1 \\
$\abseta>5.5$ & 58 & 6817 & 20.3 & 2660 & 840\,217 & 851.4 & 2718 & 847\,034 & 871.7 \\
\midrule
all & 634 & 13\,600 & 528.9 & 12\,993 & 924\,800 & 5314.6 & 13\,627 & 938\,400 & 5843.5 \\
visible ($\abseta<5.2$, no $\nu$) & 559 & 5851 & 463.4 & 9890 & 67\,472 & 4297.3 & 10\,449 & 73\,323 & 4760.7 \\
\bottomrule
\end{tabular}

\medskip
\begin{tabular}{lrrr}
\toprule
jets \code{ak4GenJets} $\to$ \code{ak4GenJetsPU} & signal (ch.~07) & signal + pileup & difference \\
\midrule
jets with $\pt>3\GeV$ & 27 & 266 & $+239$ \\
jets with $\pt>20$ / $>30$ / $>50\GeV$ & 2 / 2 / 2 & 116 / 76 / 9 & \\
signal particles inside a jet & 239 of 634 & 627 of 634 & $+388$ \\
$b$ jet: \pt{} [GeV], constituents & 117.8, 24 & 150.9, 114 (24 + 90) & $+33.1$ (pileup \pt{} 33.9, 22\,\%) \\
$\bar b$ jet: \pt{} [GeV], constituents & 115.9, 18 & 152.0, 88 (18 + 70) & $+36.1$ (pileup \pt{} 36.7, 24\,\%) \\
$\pt(\mathrm{jet})/\pt(\mathrm{quark})$, $b$ / $\bar b$ & 1.284 / 0.996 & 1.645 / 1.307 & \\
$\Delta R$(jet axis, quark), $b$ / $\bar b$ & 0.019 / 0.004 & 0.060 / 0.033 & \\
$m(b\,\mathrm{jet} + \bar b\,\mathrm{jet})$ [GeV] & 386.8 & 521.7 & $+134.9$ \\
$|\vec\pt(b) + \vec\pt(\bar b)|$ [GeV] & 6.1 & 8.6 & \\
jets without a single signal particle & -- & 87 of 266 & \\
\bottomrule
\end{tabular}
\end{table}

\begin{figure}[htbp]
\centering
\includegraphics[width=0.84\textwidth]{08_PU/out/genJetAnalyzer_08_PU.pdf}
\caption{The signal event with its 68 in-time pileup collisions, re-clustered
(\code{genJetAnalyzer\_08\_PU.pdf}): signal particles are circles, pileup particles triangles
(filled in the jet colour when clustered, grey beyond the eight leading jets; open and light grey when not); the jets above $30\GeV$ and the $b$ and
$\bar b$ jets are drawn with their circle of radius $R$ and their \pt{}, the other 190 only through their
constituents. The right-hand list gives, for the 11 leading jets, the share of their \pt{} that comes from pileup. The signal
alone is figure~\ref{fig:07_ak4} of chapter~\ref{ch:07_GenJet}.}
\label{fig:overlay}
\end{figure}

Figure~\ref{fig:overlay} shows three things. The $b$ and $\bar b$ jets kept all their signal particles, but
their axes moved by 0.04 and 0.03 in $\Delta R$ towards the pileup inside the cone and their \pt{}, their
response and their invariant mass are shifted by the full pedestal: the dijet mass goes from 387 to
$522\GeV$. Nothing tells anti-$k_{\mathrm{T}}$ which triangle belongs to the signal; it clusters what it is
given. Second, pileup makes jets of its own: 116 jets above $20\GeV$ and seven between 50 and $61\GeV$
(jets [2] to [8], 76--101\,\% pileup; jet [3] has 101 pileup constituents from 44 collisions and 8 from the
signal). These are the ``pileup jets'' that a pileup-jet identification will have to remove. Third, the
distinction is not sharp: only 87 of the 266 jets contain no signal particle, and 74 of the 76 jets above
$30\GeV$ contain at least one, because the 634 signal particles (MPI, remnants, the underlying event of
chapter~\ref{ch:04_MPI}) are spread over the whole plane and a cone anywhere catches a few. That is why 627
signal particles are now inside a jet, against 239 before: not because they moved, but because a jet above
$3\GeV$ now exists around every one of them.

\begin{plainwords}[Pileup and jets]
Spread 13\,000 soft particles evenly over the detector and draw a circle of radius 0.4 anywhere: it contains
about $35\GeV$ of them. A jet is such a circle, so every jet of the signal is some $35\GeV$ heavier than it
was, and wherever the soft particles happen to be a little denser the circle alone passes the jet threshold
and a jet appears that nobody produced. The first effect is corrected on average (the reconstruction
subtracts what a circle contains on average, $\rho\times$ area, the L1 offset of the jet energy
corrections); the second is fought by asking whether the particles in the circle come from one vertex or
from many, which is what the tracker will make possible in the next chapters. At PU 200 both numbers are
three times larger.
\end{plainwords}

\section{Open problems for Phase-2}

The list below is the state of \code{CMSSW\_15\_0\_5}, the release of this book, read in
\path{$CMSSW_RELEASE_BASE/src}, in its externals and in the global tags its \path{autoCond.py} names, not in the
development branch of CMSSW.

\paragraph{The pileup scenario.} The Phase-2 RelVals mix with \code{AVE\_200\_BX\_25ns}
(\path{Configuration/PyReleaseValidation/python/relval_steps.py}, the \code{'Run4'} branch of
\code{PUDataSets}, with \code{RelValMinBias\_14TeV} GEN-SIM as input), defined in
\path{Configuration/StandardSequences/python/Mixing.py} as \code{mix\_POISSON\_average\_cfi} with
\code{N = 200}, \code{BX = 25} and \code{B = (-3, 3)}: \code{minBunch = -3}, \code{maxBunch = 3}, seven bunch
crossings, a Poisson draw with mean 200 in each. The 2022--2025 RelVals use instead
\code{Run3\_Flat55To75\_PoissonOOTPU} (\path{SimGeneral/MixingModule/python/mix_Run3_Flat55To75_PoissonOOTPU_cfi.py}),
a flat 55--75 in-time distribution over the 16 crossings \code{minBunch = -12} to \code{maxBunch = 3} of
\path{mix_probFunction_25ns_PoissonOOTPU_cfi.py}; the 2025 profile of \path{Mixing.py}
(\path{Run3_2025_LHC_Simulation_fill12p5h_levelling9h_nb2340_xSec80mb_il2p23_cfi.py}) has the same window, a mean
of 64.3 (44.9 in 2024) and assumes ``Hadronic x-sec : 80mb'' (its header) where the library has the $78.42\mb$ of
CP5. Three consequences. The in-time pedestal of
table~\ref{tab:ledger} scales with the number of collisions: $69\GeV$ per unit area here becomes about
$200\GeV$ at PU 200, $100\GeV$ in an $R=0.4$ cone, almost what the $b$ jet of this event carries. The
out-of-time window shrinks from 16 to 7 crossings, a statement about which crossings the Phase-2 digitisers
are given (chapter~\ref{ch:11_DIGI}), not about the generator. And the library size: a PU-200 job reads on
average $200\times7 = 1400$ minimum-bias events per signal event, against $65\times16 \simeq 1040$ for the
Run-3 scenario and 193 ($71 + 68 + 54$) for the three crossings of our mixing (chapter~\ref{ch:10_SIM}); the
250 collisions of this chapter would serve one fifth of one Phase-2 event. Both release scenarios also draw the
library events at random (\code{sequential = False}: line 26 of \path{mix_POISSON_average_cfi.py}, line 42 of
\path{mix_probFunction_25ns_PoissonOOTPU_cfi.py}); only the \code{sequential = True} of this book's overlay makes
``library events 72--139 are the in-time collisions'' a meaningful statement, and a Phase-2 library smaller than
1400 events necessarily repeats collisions inside the seven crossings of one signal event.

\paragraph{The minimum-bias fragment.} \path{Configuration/Generator/python/MinBias_14TeV_pythia8_TuneCP5_cfi.py}
is the only $14\TeV$ CP5 minimum-bias fragment of the release (\path{MinBias_13TeV_pythia8_cff},
\path{SoftQCDinelastic_13TeV_pythia8_cff} and \path{SoftQCDnonDiffractive_13TeV_pythia8_cff} are CP5 at $13\TeV$,
for the generator validation; the $13.6\TeV$ \code{SoftQCD} fragments keep only events with a chosen particle; the
others are CUETP8M1, Pythia6 or heavy-ion); it is the library of the 2022--2025 RelVals (\code{RelValMinBias\_14TeV},
at $14\TeV$, without the energy override of this chapter) and of the Phase-2 ones. It imports
\path{MCTunes2017/PythiaCP5Settings_cfi.py}, whose \code{SigmaTotal:mode = 0} (``set as fixed values'',
\path{TotalCrossSections.xml} of the Pythia~8.311 external) with \code{sigmaTot = 100.309} and \code{sigmaEl =
21.89} fixes the inelastic cross section to $78.42\mb$ at every energy. The release also has a $13.6\TeV$ copy,
\path{MCTunesRun3ECM13p6TeV/PythiaCP5Settings_cfi.py}, identical but for \code{sigmaTot = 101.037} and
\code{sigmaEl = 22.08} ($78.96\mb$), imported by six python files (84 import the 2017 one), and no $14\TeV$ copy:
this chapter's library and the Phase-2 one have $78.42\mb$, the $13.6\TeV$ copy $78.96\mb$, the 2025 profile
$80\mb$. The diffractive pieces are not in the tune: \code{SigmaDiffractive:mode} keeps its default 1, the SaS
parametrisation (same Pythia file), which sets the split of table~\ref{tab:xsec} at each energy; and the
fragment's \code{crossSection = 71.39e+09} (pb) is not the $78.42\mb$ that Pythia reports.

\begingroup\sloppy
\paragraph{The vertex smearing.} The Phase-2 workflows (\code{'Run4D110'} in \path{upgradeWorkflowComponents.py})
generate a $14\TeV$ fragment with \code{-{}-beamspot DBrealisticHLLHC} (step \code{GenSimHLBeamSpot14} of
\path{relval_steps.py}), i.e.\ \path{IOMC/EventVertexGenerators/python/VtxSmearedRealisticHLLHC_cfi.py}: an
\code{HLLHCEvtVtxGenerator} that reads the payload \path{SimBeamSpotHLLHCObjects_z4p3cm_t193ns_14TeV_nominal_mc_v1}
of \code{150X\_mcRun4\_realistic\_v1} (\code{conddb list}; inserted 2024-07-04). The model is different in kind
from the Run-3 one: instead of a Gaussian in $z$ and in $t$ of the same width (\path{BetafuncEvtVtxGenerator.cc},
line 97; $3.56\cm$ and $122\ps$ for the 250 vertices of table~\ref{tab:cat}), it samples $(x,y,z,t)$ by
accept--reject from the 4D overlap density of two crossing, crab-rotated bunches (\path{HLLHCEvtVtxGenerator.cc},
lines 131--132). \code{conddb dump} of the payload gives a crossing angle of 510\,$\mu$rad, a crabbing angle of
380\,$\mu$rad at 400\,MHz, $\beta^\ast = 0.20$\,m in both planes, normalised emittances 2.5 and 2.05\,$\mu$m, a
bunch length of 0.090\,m and $7000\GeV$ per proton (the non-DB PSet of \path{VtxSmearedParameters_cfi.py}, lines
1100--1117, has 6500: ``Parameters for HL-LHC operation at 13TeV''). The time spread follows from these numbers;
the tag name gives it as ``t193ns'' (picoseconds: the 0.090\,m bunch is $300\ps$ long), and the MIP Timing
Detector (MTD), the Phase-2 layer that time-stamps charged particles with a hit resolution of $0.0370\ns$ in its
endcap (\path{RecoLocalFastTime/FTLRecProducers/python/mtdUncalibratedRecHits_cfi.py}, line 25), is what the
offline vertexing of the release uses under \code{phase2\_timing}: clusters in $z$ and $t$ (\code{DA2D\_vect},
\code{useTiming = True}, \path{RecoVertex/Configuration/python/RecoVertex_phase2_timing_cff.py}). A library
generated with the wrong beam spot cannot be repaired by the mixing, which takes every library vertex as it is; and
the overlay of this chapter, with its single bunch crossing and no time axis, cannot show that handle.
\par\endgroup
